\par\Needspace{6\baselineskip}\medskip
\hypertarget{native-118-article088}{}\pdfbookmark[2]{88.}{article-088}
\begin{center}88.\end{center}
Methodi \textit{secundae} expositio permultis relationibus novis atque elegantibus enucleandis occasionem dabit: quae quum in diversis sectionum conicarum generibus formas diversas induant, singula seorsim tractare oportebit: ab ELLIPSI initium faciemus.

\SourceLeaf{119}Respondeant duobus locis anomaliae verae $v$, $v'$ (e quibus $v$ sit tempore anterior), anomaliae excentricae $E$, $E'$, radiique vectores $r$, $r'$; porro sit $p$ semiparameter, $e=\sin\varphi$ excentricitas, $a$ semiaxis maior, $t$ tempus intra quod motus a loco primo ad secundum absolvitur; denique statuamus $v'-v=2f$, $v'+v=2F$, $E'-E=2g$, $E'+E=2G$, $a\cos\varphi=\dfrac{p}{\cos\varphi}=b$. Quibus ita factis e combinatione formularum V, VI art. 8 facile deducuntur aequationes sequentes:
\[
\begin{aligned}
\mathrm{1)}\qquad&b\sin g=\sin f\mathbin{.}\BareRoot{rr'}\\[4pt]
\mathrm{2)}\qquad&b\sin G=\sin F\mathbin{.}\BareRoot{rr'}
\end{aligned}
\]
\[
\begin{aligned}
p\cos g={}&\{\cos\tfrac12v\cos\tfrac12v'\mathbin{.}(1+e)\\[3pt]
&\quad+\sin\tfrac12v\sin\tfrac12v'\mathbin{.}(1-e)\}\BareRoot{rr'}
\end{aligned}
\]
sive
\[
\mathrm{3)}\qquad p\cos g=(\cos f+e\cos F)\BareRoot{rr'}
\]
et perinde
\[
\mathrm{4)}\qquad p\cos G=(\cos F+e\cos f)\BareRoot{rr'}.
\]
E combinatione aequationum 3, 4 porro oritur
\[
\begin{aligned}
\mathrm{5)}\qquad&\cos f\mathbin{.}\BareRoot{rr'}=(\cos g-e\cos G)a\\[4pt]
\mathrm{6)}\qquad&\cos F\mathbin{.}\BareRoot{rr'}=(\cos G-e\cos g)a.
\end{aligned}
\]
E formula III art. 8 nanciscimur
\[
\begin{aligned}
\mathrm{7)}\qquad r'-r&=2ae\sin g\sin G\\[4pt]
r'+r&=2a-2ae\cos g\cos G\\[4pt]
&=2a\sin g^2+2\cos f\cos g\mathbin{.}\BareRoot{rr'},
\end{aligned}
\]
unde
\[
\mathrm{8)}\qquad a=\frac{r+r'-2\cos f\cos g\mathbin{.}\BareRoot{rr'}}{2\sin g^2}\mathpunct{\cdot}
\]
Statuamus
\[
\mathrm{9)}\qquad\frac{\BareRoot{\dfrac{r'}r}+\BareRoot{\dfrac{r}{r'}}}{2\cos f}=1+2l,
\]
eritque
\[
\mathrm{10)}\qquad a=\frac{2(l+\sin\tfrac12g^2)\cos f\mathbin{.}\BareRoot{rr'}}{\sin g^2}
\]
nec non $\BareRoot{a}=\pm\dfrac{\BareRoot{(2(l+\sin\tfrac12g^2)\cos f\mathbin{.}\BareRoot{rr'})}}{\sin g}\mathpunct{\raisebox{.45ex}{,}}$ ubi signum superius accipere oportet vel inferius, prout $\sin g$ positivus est vel negativus. — Formula XII art. 8 \SourceLeaf{120}nobis suppeditat aequationem
\[
\begin{aligned}
\frac{kt}{a^{3/2}}&=E'-e\sin E'-E+e\sin E\\[4pt]
&=2g-2e\sin g\cos G\\[4pt]
&=2g-\sin2g+2\cos f\sin g\frac{\BareRoot{rr'}}a\mathpunct{\cdot}
\end{aligned}
\]
Quodsi iam in hac aequatione pro $a$ substituitur ipsius valor ex 10, ac brevitatis gratia ponitur
\[
\mathrm{11)}\qquad\frac{kt}{2^{3/2}\cos f^{3/2}(rr')^{3/4}}=m,
\]
prodit omnibus rite reductis
\[
\mathrm{12)}\qquad\pm m=(l+\sin\tfrac12g^2)^{1/2}+(l+\sin\tfrac12g^2)^{3/2}\left(\frac{2g-\sin2g}{\sin g^3}\right),
\]
ubi ipsi $m$ signum superius vel inferius praefigendum est, prout $\sin g$ positivus est vel negativus.

Quoties motus heliocentricus est inter $180^0$ et $360^0$, sive generalius quoties $\cos f$ est negativus, quantitas $m$ per formulam 11 determinata evaderet imaginaria, atque $l$ negativa, ad quod evitandum pro aequationibus 9, 11 in hoc casu hasce adoptabimus:
\[
\begin{aligned}
\mathrm{9^*)}\qquad&\frac{\BareRoot{\dfrac{r'}r}+\BareRoot{\dfrac{r}{r'}}}{2\cos f}=1-2L\\[8pt]
\mathrm{11^*)}\qquad&\frac{kt}{2^{3/2}(-\cos f)^{3/2}(rr')^{3/4}}=M,
\end{aligned}
\]
unde pro 10 et 12 hasce obtinebimus
\[
\begin{aligned}
\mathrm{10^*)}\qquad a&=\frac{-2(L-\sin\tfrac12g^2)\cos f\mathbin{.}\BareRoot{rr'}}{\sin g^2}\\[8pt]
\mathrm{12^*)}\qquad\pm M&=-(L-\sin\tfrac12g^2)^{1/2}+(L-\sin\tfrac12g^2)^{3/2}\left(\frac{2g-\sin2g}{\sin g^3}\right),
\end{aligned}
\]
ubi signum ambiguum eodem modo determinandum est ut ante.
